\subsection{Positive endomorphisms}

DEFINITION 6. --- Let E be a hilbertian space and \(u \in \mathcal{L}(E)\) . We say that u is positive, and write \(u \geqslant 0\) , if u is hermitian and if \(\langle x|u(x) \rangle \geqslant 0\) for all \(x \in E\) .

When K is equal to C, the relation

\[
\langle x | u (x) \rangle \geqslant 0 \quad \text { for   all } \quad x \in E
\]

implies that \(u\) is hermitian (V, p.~2, Remark), hence positive.

Let \(\mathcal{L}_{+}(\mathrm{E})\) denote the set of all positive elements of \(\mathcal{L}(\mathrm{E})\) ; this is a proper pointed convex cone in the real vector space \(\mathcal{L}(\mathrm{E})_{[\mathbb{R}]}\) underlying \(\mathcal{L}(\mathrm{E})\) . In order that \(u\) is positive, it is necessary and sufficient that the sesquilinear form \(\Phi_u\) on \(\mathrm{E} \times \mathrm{E}\) associated with \(u\) is positive hermitian. Given \(u\) and \(v\) in \(\mathcal{L}(\mathrm{E})\) , the relation \(u - v \geqslant 0\) can also be written as \(u \geqslant v\) or \(v \leqslant u\) ; this is an order relation on \(\mathcal{L}(\mathrm{E})_{[\mathbb{R}]}\) compatible with its real vector space structure.

PROPOSITION 12. --- Let u be a hermitian (resp. positive) element of \(\mathcal{L}(\mathrm{E})\) and let v be a continuous linear mapping from E into a hilbertian space F. Then vuv* is a hermitian (resp. positive) element of \(\mathcal{L}(\mathrm{F})\) .

For, we have \((vuv^{*})^{*} = v^{**}u^{*}v^{*} = vuv^{*}\) . On the other hand, if \(u \geqslant 0\) , we have

\[
\langle y | v u v ^ {*} (y) \rangle = \langle v ^ {*} (y) | u (v ^ {*} (y)) \rangle \geqslant 0
\]

for all \(y \in \mathbf{F}\) , hence \(vuv^{*} \geqslant 0\) .

Prop. 12 shows, in particular, that \(vv^*\) is positive for all \(v \in \mathcal{L}(\mathrm{E};\mathrm{F})\) . Since, in particular, an orthoprojector \(p\) satisfies \(p = p^2 = pp^*\) , it is positive.

Remarks. --- 1) For every hermitian \(u\) in \(\mathcal{L}(\mathrm{E})\) , put \(m(u) = \inf_{||x|| = 1} \langle x|u(x) \rangle\) , \(\mathbf{M}(u) = \sup_{||x|| = 1} \langle x|u(x) \rangle\) . If E is not just 0, \(m(u)\) and \(\mathbf{M}(u)\) are finite; moreover, \(\mathbf{M}(u)\) is the smallest real number \(\lambda\) such that \(u \leqslant \lambda\) . \(1_{\mathrm{E}}\) and \(m(u)\) the largest real number \(\mu\) such that \(u \geqslant \mu\) . \(1_{\mathrm{E}}\) . Clearly we have \(m(-u) = -\mathbf{M}(u)\) and \(\mathbf{M}(-u) = -m(u)\) . It is clear that

\[
\sup \left(| m (u) |, | \mathrm{M} (u) |\right) = \sup _ {\| x \| = 1} \left| \langle x | u (x) \rangle \right|
\]

and prop. 9 (V, p.~44) implies (for E ≠ \{0\}) that

\[
\| u \| = \sup \left(| m (u) |, | \mathbf {M} (u) |\right).\tag{17}
\]

* For another proof of this formula when K is C, see prop. 14 of TS, I, § 6, No.~8. * 2) Let M and N be two closed vector subspaces of E, and \(p_{\mathbf{M}}\) (resp. \(p_{\mathbf{N}}\) ) the orthoprojector from E onto M (resp. N). Then \(\mathbf{M} \subset \mathbf{N}\) if and only if \(p_{\mathbf{M}} \leqslant p_{\mathbf{N}}\) . For, we have \(p_{\mathbf{M}}^{*}p_{\mathbf{M}} = p_{\mathbf{M}}\) , hence

\[
\left\| p _ {\mathbf {M}} (x) \right\| ^ {2} = \langle p _ {\mathbf {M}} (x) | p _ {\mathbf {M}} (x) \rangle = \langle x | p _ {\mathbf {M}} ^ {*} p _ {\mathbf {M}} (x) \rangle = \langle x | p _ {\mathbf {M}} (x) \rangle
\]

for all \(x \in E\) . The relation \(p_{M} \leqslant p_{N}\) is therefore equivalent to « \(\|p_{M}(x)\| \leqslant \|p_{N}(x)\|\) for all \(x \in E\) ». If \(M \subset N\) , we have \(p_{M} = p_{M} p_{N}\) , hence \(\|p_{M}(x)\| \leqslant \|p_{N}(x)\|\) since \(\|p_{M}\| \leqslant 1\) . Conversely, if \(\|p_{M}(x)\| \leqslant \|p_{N}(x)\|\) for all \(x \in E\) , the kernel of \(p_{M}\) contains the kernel of \(p_{N}\) , that is, that \(M^{\circ} \supset N^{\circ}\) , which implies that \(M \subset N\) .

PROPOSITION 13. --- Let \(\mathcal{H}(\mathrm{E})\) be the set of all continuous hermitian endomorphisms of the hilbertian space E. Let \(\mathcal{F}\) be a non-empty, directed increasing and bounded subset of \(\mathcal{H}(\mathrm{E})\) .

\begin{enumerate}
\def\labelenumi{(\roman{enumi})}
\tightlist
\item
  The set \(\mathcal{F}\) has an upper bound \(u_0\) in \(\mathcal{H}(\mathrm{E})\) ; we have
\end{enumerate}

\[
\langle x | u _ {0} (x) \rangle = \sup _ {u \in \mathcal {F}} \langle x | u (x) \rangle \quad f o r a l l \quad x \in E.\tag{18}
\]

\begin{enumerate}
\def\labelenumi{(\roman{enumi})}
\setcounter{enumi}{1}
\tightlist
\item
  The filter of sections of \(\mathcal{F}\) converges to \(u_0\) in the space \(\mathcal{L}(\mathrm{E})\) endowed with the topology of simple convergence.
\end{enumerate}

Let \(\Sigma\) be the filter of sections of \(\mathcal{F}\) ; for every \(u \in \mathcal{H}(\mathrm{E})\) , let \(\Phi_u\) be the continuous hermitian form on \(\mathbf{E}\) defined by

\[
\Phi_ {u} (x, y) = \langle x | u (y) \rangle .
\]

Let

\[
\Psi_ {u} (x) = \Phi_ {u} (x, x)
\]

for \(u \in \mathcal{H}(\mathrm{E})\) and \(x \in \mathrm{E}\) . By the polarization formulas (V, p.~2), we have

\[
4 \Phi_ {u} (x, y) = \Psi_ {u} (x + y) - \Psi_ {u} (x - y) \quad \text { if } \quad \mathbf {K} = \mathbf {R}\tag{19}
\]

\[
4 \Phi_ {u} (x, y) = \Psi_ {u} (x + y) - \Psi_ {u} (x - y) - i \Psi_ {u} (x + i y) + i \Psi_ {u} (x - i y) \quad \text { if } \quad \mathbf {K} = \mathbf {C}.\tag{20}
\]

For every \(x \in \mathbf{E}\) , the mapping \(u \mapsto \Psi_u(x)\) from into \(\mathbf{R}\) is increasing and bounded, hence has a limit with respect to \(\Sigma\) . By the preceding formulas, the limit

\[
\lim _ {u, \Sigma} \Phi_ {u} (x, y) = \Phi (x, y)
\]

exists for every pair \((x, y)\) of elements of E. It is clear that \(\Phi\) is a hermitian form on E. If \(v_{1} \in F\) and \(v_{2}\) is a bound of F, the hermitian forms \(f_{1} = \Phi - \Phi_{v_{1}}\) and \(f_{2} = \Phi_{v_{2}} - \Phi\) are positive; there exists a real number \(M \geqslant 0\) such that

\[
f _ {1} (x, x) + f _ {2} (x, x) = \Phi_ {v _ {2} - v _ {1}} (x, x) \leqslant \mathbf {M} \| x \| ^ {2},
\]

hence

\[
f _ {1} (x, x) \leqslant \mathbf {M} \| x \| ^ {2}, f _ {2} (x, x) \leqslant \mathbf {M} \| x \| ^ {2} (x \in \mathrm{E});
\]

consequently the semi-norms \(x \mapsto f_i(x, x)^{1/2}\) are continuous on E. Since

\[
f _ {2} - f _ {1} = \Phi_ {v _ {2}} + \Phi_ {v _ {1}} - 2 \Phi ,
\]

we conclude that \(x \mapsto \Phi(x, x)\) is a continuous function on E, and by formulas (19) and (20), that \(\Phi\) is continuous on \(E \times E\) . Therefore there exists (V, p.~16, cor. 2) an element \(u_{0}\) of \(\mathcal{H}(E)\) such that \(\Phi = \Phi_{u_{0}}\) . Formula (18) is evidently satisfied, hence \(u_{0}\) is the upper bound of F in \(\mathcal{H}(E)\) . This proves (i).

We have, by construction

\[
\lim _ {u, \Sigma} \langle x | (u _ {0} - u) (x) \rangle = 0 \quad \text { for   all } \quad x \in E.\tag{21}
\]

Let \(v_{1} \in \mathcal{F}\) ; given a \(u \in \mathcal{F}\) such that \(u \geqslant v_{1}\) , let \(v = u_{0} - u\) . If we apply the Cauchy-Schwarz inequality to the positive hermitian form \(\Phi_v\) on E, we get

\[
\begin{array}{r l} \left\| v (x) \right\| ^ {4} & = \left| \Phi_ {v} (v (x), x) \right| ^ {2} \leqslant \Phi_ {v} (v (x), v (x)). \Phi_ {v} (x, x) \\ & = \langle v (x) | v ^ {2} (x) \rangle \langle x | v (x) \rangle \leqslant \| v \| ^ {3} \| x \| ^ {2} \langle x | v (x) \rangle \\ & \leqslant \| u _ {0} - v _ {1} \| ^ {3} \| x \| ^ {2} \langle x | v (x) \rangle , \end{array}
\]

since \(\| v\| \leqslant \| u_0 - v_1\|\) by V, p.~44, prop. 9. Then by (21) we get \(\lim_{u,\Sigma}\left\| (u_0 - u)(x)\right\| = 0\) for all \(x\in \mathrm{E}\) ; which proves assertion (ii). Q.E.D.

In particular, prop. 13 can be applied to the case of an increasing and bounded sequence \((u_n)_{n\in \mathbf{N}}\) of elements of \(\mathcal{H}(\mathrm{E})\) . Then there exists an element \(v\) of \(\mathcal{H}(\mathrm{E})\) characterized by

\[
\langle x | v (x) \rangle = \lim _ {n \rightarrow \infty} \langle x | u _ {n} (x) \rangle = \sup _ {n \in \mathbf {N}} \langle x | u _ {n} (x) \rangle \quad (x \in E),
\]

and we have \(v(x) = \lim_{n\to \infty}u_n(x)\) for all \(x\in E\) . Moreover, \(v\) is the upper bound of the set of the \(u_{n}\) in \(\mathcal{H}(E)\) .

